\chapter{Materiel et fondements theoriques}

\section{Presentation du materiel}

\subsection{Presentation de la camera DAVIS346 et fiche technique}
La camera DAVIS346 est un capteur neuromorphique hybride qui produit en
parallele un flux d'evenements (DVS) et des images APS classiques sur un meme
capteur. Cette architecture est interessante pour notre projet, car elle permet
de conserver une information temporelle tres fine tout en gardant une reference
image exploitable pour le debug et la visualisation \cite{davis346_datasheet}.

Les caracteristiques techniques directement utiles pour la chaine de perception
du projet sont resumees dans le tableau \ref{tab:davis_specs}.

\begin{table}[H]
\centering
\caption{Fiche technique synthetique de la DAVIS346}
\label{tab:davis_specs}
\begin{tabular}{ll}
\toprule
Parametre & Valeur \\
\midrule
Resolution capteur & 346 x 260 \\
Sorties & Evenements DVS + images APS concurrentes \\
Plage dynamique evenementielle & Jusqu'a 120 dB \\
Resolution temporelle & 1 us \\
Latence evenementielle & Sous 1 us (nominale) \\
Debit maximal & Jusqu'a 12 MEPS \\
IMU integree & 6 axes (gyro + accelerometre), jusqu'a 8 kHz \\
Alimentation & $<$ 180 mA a 5 V \\
Connectique principale & USB 3.0 Micro-B (verrouillable) \\
Monture optique & CS-mount (adaptateur C-mount possible) \\
\bottomrule
\end{tabular}
\end{table}

Le document technique signale egalement des limites pratiques a connaitre lors
des essais: salves d'evenements possibles en mode APS Global Shutter, schema de
balayage DVS sous forte charge, et performances image APS inferieures a des
cameras classiques. Ces limites justifient un reglage fin des biais capteur et
une calibration experimentale adaptee au contexte robotique
\cite{davis346_datasheet}.

\subsection{Presentation du robot LIMO}
Le robot LIMO d'AgileX est presente comme une plateforme ROS polyvalente avec
quatre modes de locomotion (differentiel, Ackermann, chenilles et mecanum), ce
qui en fait un support pertinent pour des tests de navigation dans des scenes
heterogenes \cite{limo_ros2_doc,limo_agilex_video}.

La documentation technique indique en particulier:
\begin{itemize}
\item dimensions compactes (322 x 220 x 251 mm) et charge utile de 4 kg;
\item vitesse maximale sans charge de l'ordre de 1 m/s;
\item calcul embarque base sur NUC sous Ubuntu 22.04;
\item capteurs integres utiles a la navigation (LiDAR, camera de profondeur,
IMU).
\end{itemize}

La pile logicielle officielle propose des launches ROS 2 pour le demarrage de
la base, la cartographie, puis la navigation Nav2 suivant le mode de conduite
retendu \cite{limo_ros2_doc}. Cette base materielle et logicielle est
compatible avec notre integration de la perception evenementielle.

\section{Concepts theoriques mobilises}

\subsection{Camera a evenements: notion, bref historique, avantages et limites}
Une camera a evenements est un capteur asynchrone: chaque pixel declenche un
evenement uniquement lors d'une variation de luminance, au lieu d'attendre une
image globale echantillonnee a frequence fixe. La revue de reference rappelle
que ce paradigme s'inscrit dans le champ plus large de l'ingenierie
neuromorphique et qu'il s'est structure au cours des dernieres decennies avec
des familles de capteurs comme DVS ou ATIS \cite{novoa2024review}.

Les avantages majeurs pour la robotique mobile sont:
\begin{itemize}
\item forte plage dynamique (ordre de 120 dB) utile en eclairage difficile;
\item tres faible latence de perception;
\item reduction de la redondance et de la consommation grace au codage
asynchrone des changements \cite{novoa2024review}.
\end{itemize}

Les limites relevees dans la meme revue sont aussi importantes:
\begin{itemize}
\item flux sparse plus difficile a exploiter que des images denses;
\item capteurs encore couteux;
\item resolution et interfaces encore contraintes sur certains systemes
commerciaux \cite{novoa2024review}.
\end{itemize}

\subsection{Detection d'obstacles (IMO): panorama rapide des familles}
La detection d'objets en mouvement independant (IMO) vise a separer les
mouvements lies a l'ego-mouvement du robot de ceux des objets reels de la
scene. Dans un cadre evenementiel, la revue distingue trois familles
methodologiques principales \cite{novoa2024review}:
\begin{itemize}
\item approches de compensation de mouvement (soustraction du mouvement global,
souvent via flot optique);
\item approches residuelles iteratives combinant ego-mouvement et segmentation
de residus;
\item approches apprentissage (supervise ou non supervise) pour segmenter
directement les IMO.
\end{itemize}

Ce panorama montre que les approches a base de compensation restent attractives
pour des contraintes de temps reel embarque, ce qui motive le choix methodologique
du projet.

\subsection{Navigation autonome et algorithme utilise dans le projet}
La documentation LIMO rappelle le schema classique de navigation mobile:
localisation sur carte + planification globale + planification locale
\cite{limo_ros2_doc}. En pratique, l'execution de la navigation est faite via
Nav2 avec un lancement dedie au mode differentiel
(\texttt{limo\_nav2\_diff.launch.py}) \cite{limo_ros2_doc}.

Dans notre configuration, le planificateur global est de type NavFn
(equivalent Dijkstra quand l'option A* est inactive), le controle local repose
sur DWB, et la localisation est assuree par AMCL. Cette architecture sert de
base pour injecter ensuite les obstacles dynamiques issus de la perception
evenementielle.
